% ═══════════════════════════════════════════════════════════════════
%  INFORME — Tarea IC RNA 2026-1   (v2, revisado)
%  Paper: Widrow & Lehr (1990)
%  "30 Years of Adaptive Neural Networks:
%   Perceptron, Madaline, and Backpropagation"
% ═══════════════════════════════════════════════════════════════════

\documentclass[11pt,a4paper]{article}

% ── Paquetes ──────────────────────────────────────────────────────
\usepackage[utf8]{inputenc}
\usepackage[T1]{fontenc}
\usepackage[spanish,es-noshorthands]{babel}
\usepackage{amsmath,amssymb,amsfonts}
\usepackage{graphicx}
\usepackage[margin=2.0cm]{geometry}
\usepackage{float}
\usepackage{booktabs}
\usepackage{enumitem}
\usepackage{caption}
\captionsetup{font=small,skip=4pt}
\emergencystretch=2.5em
\usepackage{tikz}
\usetikzlibrary{arrows.meta}
\usepackage[colorlinks=true,linkcolor=blue,citecolor=blue,urlcolor=blue]{hyperref}

% ── Macros para figuras y tablas (compilacion robusta) ───────────
\newcommand{\figura}[4][0.78]{%
\begin{figure}[htbp]\centering
\IfFileExists{../../notebook/v1/figs/#2.png}{%
\includegraphics[width=#1\textwidth]{../../notebook/v1/figs/#2.png}}{%
\fbox{\parbox[c][2.6cm][c]{0.8\textwidth}{\centering\texttt{figs/\detokenize{#2}.png}\\(pendiente: ejecutar el notebook)}}}
\caption{#3}\label{fig:#4}
\end{figure}}

\newcommand{\tabla}[3]{%
\begin{table}[htbp]\centering\small
\IfFileExists{../../notebook/v1/tabs/#1.tex}{%
\input{../../notebook/v1/tabs/#1.tex}}{%
\fbox{\parbox[c][1.6cm][c]{0.7\textwidth}{\centering\texttt{tabs/\detokenize{#1}.tex}\\(pendiente)}}}
\caption{#2}\label{tab:#3}
\end{table}}

% ── Comandos personalizados ──────────────────────────────────────
\newcommand{\R}{\mathbb{R}}
\newcommand{\E}{\mathbb{E}}
\newcommand{\norm}[1]{\lVert#1\rVert}
\newcommand{\aLMS}{$\alpha$-LMS}
\newcommand{\muLMS}{$\mu$-LMS}
\newcommand{\bp}{\textit{backpropagation}}
\newcommand{\mri}{\textit{Madaline Rule I}}
\newcommand{\mrii}{\textit{Madaline Rule II}}
\newcommand{\mriii}{\textit{Madaline Rule III}}

% ── Metadatos ─────────────────────────────────────────────────────
\title{\textbf{Algoritmos de Aprendizaje en Redes Neuronales Feedforward}\\[4pt]
\large Un análisis del artículo de Widrow \& Lehr (1990):\\
\textit{30 Years of Adaptive Neural Networks: Perceptron, Madaline, and Backpropagation}}
\author{Estudiante\\Universidad de Santiago de Chile\\Magíster en Ingeniería Informática}
\date{Julio 2026}

\begin{document}
\maketitle

\begin{abstract}
Este informe presenta un análisis del artículo de revisión de Widrow \& Lehr (1990)
sobre 30 años de algoritmos de aprendizaje supervisado en redes neuronales artificiales (RNA)
feedforward. El marco teórico cubre la familia completa de algoritmos del artículo, organizada
según los dos ejes del propio paper: reglas de \textbf{corrección de error} (\aLMS, Perceptrón,
Mays, \mri, \mrii) y reglas de \textbf{descenso más pronunciado} (\muLMS, \bp, \mriii). De ellos
se implementaron desde cero en NumPy puro los seis algoritmos experimentables
—Perceptrón, \aLMS, \muLMS, MLP con \bp, \mrii~y \mriii— y se ejecutaron siete experimentos
para verificar las hipótesis fundamentales del principio de mínima perturbación.
Los resultados confirman que: (1) la capacidad estadística $C_s \approx 2N_w$ de Cover se verifica
empíricamente, (2) el Perceptrón converge en tiempo finito solo en datos separables mientras \muLMS
se aproxima al óptimo de Wiener, (3) el MLP con \bp~resuelve XOR (6/20, 2--2--1) y alcanza 91.1\% en
Iris, (4) \mriii~es matemáticamente equivalente a \bp~con error relativo de $5.05 \times 10^{-8}$,
a un costo computacional que crece aproximadamente como $\mathcal{O}(N_{\text{Adalines}})$. Se discute
el principio de mínima perturbación como marco unificador del aprendizaje supervisado en RNA.
\end{abstract}


% ═══════════════════════════════════════════════════════════════════
%  1. INTRODUCCIÓN
% ═══════════════════════════════════════════════════════════════════
\section{Introducción}

\subsection{Problema abordado}

El artículo de Widrow \& Lehr (1990) \cite{widrow1990} aborda el problema fundamental del
\textbf{aprendizaje supervisado en redes neuronales artificiales feedforward}: ¿cómo adaptar los
pesos sinápticos de una red de elementos computacionales simples para que aproximen una función
de clasificación o regresión a partir de ejemplos? El artículo revisa 30 años de desarrollo de
algoritmos de entrenamiento —desde el Perceptrón de Rosenblatt (1958) hasta la retropropagación
del error de Rumelhart, Hinton \& Williams (1986)— y demuestra que todos ellos comparten un
principio unificador: el \textbf{principio de mínima perturbación} (\textit{minimal disturbance
principle}), que establece que durante el entrenamiento es recomendable inyectar nueva información
en la red de manera que perturbe la información ya almacenada en la menor medida posible
\cite{widrow1990}.

El artículo es una revisión (\textit{survey}) y tutorial, no un estudio experimental sobre un
dataset único. Su contribución es \textbf{conceptual e histórica}: traza la evolución desde las
primeras reglas de aprendizaje (1960) hasta las redes multicapa con \bp, y demuestra conexiones
matemáticas profundas entre algoritmos desarrollados independientemente. En particular, establece
que \mriii~(Andes, 1988) es matemáticamente equivalente a \bp, un hallazgo realizado por Widrow
y sus estudiantes \cite{widrow1990}.

\subsection{Reseña histórica}

La Sección I del artículo sitúa la contribución dentro de tres décadas de desarrollo. Los hitos
del linaje que el paper sigue en detalle son: la \textbf{regla del Perceptrón} (Rosenblatt) y el
algoritmo \textbf{LMS} (Widrow \& Hoff), ambos publicados en 1960; la \textbf{Madaline Rule I}
(mediados de los 60), primera regla popular para redes con múltiples elementos adaptativos, con
aplicaciones en voz, predicción del clima y reconocimiento de patrones; la \textbf{retropropagación}
formulada por Werbos (1971--1974), ignorada inicialmente, redescubierta por Parker (1982) y
popularizada por Rumelhart, Hinton \& Williams (1986); la \textbf{Madaline Rule II} (Widrow, Winter
\& Baxter, 1987) para adaptar redes multicapa con cuantizadores duros; y la \textbf{Madaline Rule
III} (Andes, 1988). En paralelo, el artículo reconoce ramas que no desarrolla en profundidad
—la \textit{Learning Matrix} de Steinbuch, la teoría ART de Grossberg, los mapas de Kohonen, los
modelos de Hopfield, la máquina de Boltzmann y el Cognitron/Neocognitron de Fukushima—, delimitando
que su foco son los algoritmos nacidos en Stanford y la retropropagación \cite{widrow1990}.

\subsection{¿Por qué redes neuronales artificiales?}

La motivación para usar RNA en problemas de clasificación y regresión se fundamenta en tres
argumentos presentados en las Secciones I--II del artículo \cite{widrow1990}:

\begin{enumerate}[label=\textbf{(\roman*)},leftmargin=*]
    \item \textbf{Paralelismo masivo y robustez:} los sistemas biológicos resuelven tareas de
    reconocimiento de patrones y procesamiento de voz e imágenes con una facilidad que las
    computadoras von Neumann no igualan. Las RNA emulan este paralelismo mediante muchos elementos
    simples interconectados que procesan información simultáneamente.
    \item \textbf{Aprendizaje por ejemplos:} a diferencia de los sistemas algorítmicos
    tradicionales, las RNA no requieren programación explícita de reglas; aprenden el mapeo
    entrada--salida directamente de los datos mediante algoritmos iterativos.
    \item \textbf{Generalización:} una RNA correctamente entrenada clasifica patrones no vistos
    durante el entrenamiento. La capacidad de generalización está inversamente relacionada con la
    capacidad de almacenamiento ($C_s \approx 2N_w$, Cover, 1964 \cite{cover1965}), un
    \textit{trade-off} fundamental que el artículo explora en profundidad.
\end{enumerate}

Aplicaciones históricas que demuestran la viabilidad práctica incluyen NETtalk (Sejnowski \&
Rosenberg, 1987), una red de 80--26 Adalines que aprende a pronunciar texto en inglés; el
\textit{truck backer-upper} de Nguyen \& Widrow (1989), donde una red de 26 Adalines aprende a
retroceder un camión con acoplado; y los ecualizadores adaptativos basados en LMS para módems
de alta velocidad, la primera aplicación comercial masiva de las RNA \cite{widrow1990}.

\subsection{Hipótesis}

Se plantean cuatro hipótesis que guían la experimentación:

\begin{description}
    \item[H1 -- Capacidad de Cover:] la capacidad estadística de un Adaline es $C_s \approx 2N_w$
    y la determinística $C_d = N_w$, verificable mediante simulación Monte Carlo con programación
    lineal.
    \item[H2 -- Convergencia finita del Perceptrón:] el Perceptrón converge en tiempo finito si
    los datos son linealmente separables, pero oscila y diverge en norma si no lo son. Por
    contraste, \aLMS~y \muLMS~son estables en ambos casos, con \muLMS~convergiendo al óptimo de
    Wiener $W^* = R^{-1}P$.
    \item[H3 -- Poder no lineal y óptimos locales:] XOR es irresoluble por un solo Adaline pero
    resoluble por \mrii~y MLP con \bp, aunque ambos algoritmos están sujetos a óptimos locales
    en la superficie de error. El remedio de ruido esporádico propuesto en la Sección VII-E del
    paper \cite{widrow1990} mitiga parcialmente este problema.
    \item[H4 -- Equivalencia MRIII $\equiv$ Backpropagation:] \mriii~converge a \bp~cuando la
    perturbación $\Delta s \to 0$, pero su costo computacional en implementación digital es
    proporcionalmente mayor (requiere $1 + N_{\text{Adalines}}$ pasadas hacia adelante por
    presentación versus 1 adelante $+$ 1 atrás de \bp).
\end{description}

% ═══════════════════════════════════════════════════════════════════
%  2. MARCO TEÓRICO  (reorganizado según los dos ejes de la Fig. 33)
% ═══════════════════════════════════════════════════════════════════
\section{Marco Teórico}

El artículo organiza todos sus algoritmos según dos ejes (Fig.~\ref{fig:taxonomia}): el
\textbf{mecanismo de adaptación} —corrección de error frente a descenso más pronunciado— y la
\textbf{no linealidad} del elemento —lineal, sigmoide o signum—, para uno o múltiples elementos.
Este marco teórico sigue esa misma estructura: primero los conceptos comunes
(§\ref{ssec:fundamentos}--\ref{ssec:mdp}), luego la familia de corrección de error
(§\ref{ssec:ec}) y la de descenso más pronunciado (§\ref{ssec:sd}), cerrando con la taxonomía
(§\ref{ssec:taxonomia}).

\subsection{Conceptos fundamentales}\label{ssec:fundamentos}

\subsubsection{El combinador lineal adaptativo y el Adaline}

El bloque constructivo fundamental de las RNA feedforward es el \textbf{combinador lineal
adaptativo} (Fig.~1 del paper), que calcula el producto interno $s_k = X_k^T W_k$, donde
$X_k = [x_0=+1, x_{1k}, \ldots, x_{nk}]^T$ es el vector de entrada aumentado con bias y
$W_k \in \R^{n+1}$ el vector de pesos. Cuando su salida lineal se hace pasar por un
\textbf{cuantizador} de límite duro (\textit{signum}) se obtiene el \textbf{Adaptive Linear
Element} (Adaline, Fig.~2), un elemento lógico de umbral adaptativo con salida binaria
$y_k = \mathrm{sgn}(s_k) \in \{\pm1\}$. El peso de bias $w_{0k}$ controla el umbral. En su
variante diferenciable, el signum se reemplaza por una sigmoide como la tangente hiperbólica
$y_k = \tanh(s_k)$, que habilita las reglas de descenso por gradiente \cite{widrow1990}.

\subsubsection{Separabilidad lineal y capacidad de patrones}

La condición crítica $s = 0$ define un \textbf{hiperplano separador} en el espacio de patrones.
Con $n$ entradas binarias, un Adaline solo puede implementar las funciones lógicas
\textbf{linealmente separables}: 14 de 16 posibles para $n=2$, excluyendo XOR/XNOR
\cite{widrow1990}. Cover (1964) \cite{cover1965} demostró que la probabilidad de que $N_p$
patrones aleatorios sean linealmente separables por un Adaline de $N_w$ pesos es:

\begin{equation}
P(N_p, N_w) =
\begin{cases}
1 & \text{si } N_p \leq N_w \\[4pt]
2^{1-N_p} \displaystyle\sum_{i=0}^{N_w-1} \binom{N_p-1}{i} & \text{si } N_p > N_w
\end{cases}
\label{eq:cover}
\end{equation}

De aquí se derivan la \textbf{capacidad determinística} $C_d = N_w$ (garantía de solución) y la
\textbf{capacidad estadística} $C_s \approx 2N_w$ (promedio de patrones absorbibles). Para redes
multicapa, Baum acotó $N_w/N_y - K_1 \leq C_d \leq (N_w/N_y)\log_2(N_w/N_y) + K_2$
\cite{widrow1990}. La regla práctica de generalización es $N_p \gg N_w/N_y$: el número de patrones
de entrenamiento debe exceder la relación pesos/salidas para que la red generalice en lugar de
memorizar.

\subsubsection{Clasificadores no lineales: el preprocesador polinomial}

La primera vía del artículo hacia fronteras no lineales sin usar varias capas es un
\textbf{preprocesador polinomial fijo} conectado a un único Adaline adaptativo (Fig.~6 del paper).
Para dos entradas, el preprocesador expande el vector a los términos $[x_1^2, x_1, x_1 x_2, x_2,
x_2^2]$ y la condición de umbral se vuelve:

\begin{equation}
s = w_0 + x_1 w_1 + x_1^2 w_{11} + x_1 x_2 w_{12} + x_2^2 w_{22} + x_2 w_2 = 0.
\label{eq:polinomial}
\end{equation}

Con una elección adecuada de pesos, la frontera separadora en el espacio de patrones puede ser una
curva cerrada —por ejemplo, una \textbf{elipse}— que sí resuelve el XNOR (Fig.~7). Especht
desarrolló para esto el \textit{polynomial discriminant method} (PDM), un procedimiento no
iterativo de una sola pasada. Tras el entrenamiento por gradiente, los pesos del Adaline aproximan
los coeficientes de una expansión de Taylor multidimensional de la función de respuesta deseada
\cite{widrow1990}. La elección del preprocesador determina la generalización; la experiencia indica
que, salvo que las no linealidades se ajusten con cuidado al problema, suele obtenerse mejor
generalización con redes de más de una capa adaptativa.

\subsubsection{Redes feedforward y las señales de error ocultas}

La alternativa general es la \textbf{red feedforward multicapa} totalmente conectada (Fig.~11),
donde todas las capas son adaptativas. La adaptación de la capa de salida es directa —el error se
computa comparando con la respuesta deseada—, pero la dificultad fundamental es obtener
``\textbf{señales de error}'' para las Adalines de las capas ocultas. \bp~y \mriii~son
precisamente los dos mecanismos que el artículo propone para establecer estas señales
\cite{widrow1990}.

\subsection{El principio de mínima perturbación}\label{ssec:mdp}

Todos los algoritmos del artículo obedecen un principio común (Sección III): \emph{adaptar para
reducir el error de salida del patrón actual, con la mínima perturbación a las respuestas ya
aprendidas}. Geométricamente, esto se traduce en que el cambio de pesos $\Delta W_k$ sea
\textbf{colineal} con el vector de entrada $X_k$, logrando la corrección deseada con el cambio de
menor magnitud posible. Este principio, intuitivo, fue lo que condujo al descubrimiento del
algoritmo LMS y de las reglas Madaline \cite{widrow1990}.

\subsection{Reglas de corrección de error}\label{ssec:ec}

Las reglas de corrección de error alteran los pesos para corregir una proporción del error en la
respuesta al patrón presente. Se aplican primero a un solo Adaline y luego a redes.

\subsubsection{$\alpha$-LMS — regla delta de Widrow-Hoff (Eq.~10)}

\begin{equation}
W_{k+1} = W_k + \alpha\,\frac{\varepsilon_k}{\norm{X_k}^2}\,X_k,\quad \varepsilon_k = d_k - X_k^T W_k
\label{eq:alpha_lms}
\end{equation}

Única regla lineal de corrección de error: el cambio es proporcional al error \emph{lineal} y
colineal con $X_k$. Es auto-normalizante (la elección de $\alpha$ no depende de la magnitud de las
entradas). Rango de estabilidad $0 < \alpha < 2$; práctico $0.1 < \alpha < 1.0$ \cite{widrow1990}.

\subsubsection{Regla del Perceptrón de Rosenblatt (Eq.~18)}

\begin{equation}
W_{k+1} = W_k + \alpha\,\frac{\tilde{\varepsilon}_k}{2}\,X_k,\quad \tilde{\varepsilon}_k = d_k - y_k \in \{-2,0,+2\}
\label{eq:perceptron}
\end{equation}

Regla no lineal: usa el error del \emph{cuantizador} $\tilde{\varepsilon}_k$ y solo adapta si
$y_k \neq d_k$. El \textbf{teorema de convergencia del Perceptrón} garantiza separación en pasos
finitos si los datos son linealmente separables; si no lo son, el vector de pesos ``gravita hacia
cero'' y el algoritmo continúa indefinidamente \cite{widrow1990,rosenblatt1958}.

\subsubsection{Reglas de Mays: la zona muerta}

Mays (1963) introdujo una \textbf{zona muerta} $\gamma$ alrededor de la salida lineal $s_k$. En su
regla de \textit{adaptación por incremento} \cite{widrow1990}:

\begin{equation}
W_{k+1} =
\begin{cases}
W_k + \alpha\,\tilde{\varepsilon}_k\,\dfrac{X_k}{2\norm{X_k}^2} & \text{si } |s_k| \geq \gamma \\[8pt]
W_k + \alpha\,d_k\,\dfrac{X_k}{\norm{X_k}^2} & \text{si } |s_k| < \gamma
\end{cases}
\label{eq:mays}
\end{equation}

Si los patrones son linealmente separables, la regla converge en pasos finitos. La ventaja aparece
en datos \emph{no} separables: una zona muerta suficientemente grande tiende a alejar el vector de
pesos de cero cuando existe una solución razonablemente buena, por lo que Mays se desempeña mucho
mejor que el Perceptrón en ese régimen. Los casos límite recuperan reglas conocidas: con
$\gamma \to 0$ se obtiene el Perceptrón normalizado, y la variante de \textit{relajación
modificada} con $\gamma \to \infty$ se reduce a \aLMS~\cite{widrow1990}.

\subsubsection{Madaline Rule I (MRI)}

MRI (Widrow \& Hoff, principios de los 60) es una de las primeras redes en capas entrenables:
una capa de Adalines con cuantizador signum cuyas salidas alimentan un \textbf{elemento lógico fijo}
de segunda capa (AND, OR o mayoría MAJ; Figs.~8--10 y 15 del paper). Como la lógica de salida es
conocida, es posible determinar qué Adalines de la primera capa deben adaptarse para corregir un
error. La regla adapta la(s) Adaline(s) cuya salida lineal $|s_k|$ está \emph{más cerca de cero}
—las que requieren el menor cambio de pesos para invertir su respuesta— invirtiendo las suficientes
para que la lógica de salida coincida con la deseada. La adaptación puede ser \emph{absoluta}
(corrección agresiva, aprendizaje ``rápido'') o por el incremento de \aLMS~(``lento''). Un
\textit{job assigner} reparte la responsabilidad entre Adalines (\textit{load-sharing}: ``asignar
la responsabilidad a quien pueda asumirla más fácilmente''), y la presentación de patrones debe ser
aleatoria para evitar ciclos \cite{widrow1990}. Como MRII, MRI puede estancarse en óptimos locales.

\subsubsection{Madaline Rule II (MRII)}

\mrii~(Widrow, Winter \& Baxter, 1987) \cite{widrow1990} extiende la adaptación a \emph{todas} las
capas de una red de Adalines con cuantizador signum (Fig.~16). Para cada patrón erróneo:
(1)~se ordenan las Adalines por $|s|$ creciente; (2)~se invierte tentativamente (\textit{flip}) la
salida de una Adalina; si reduce el error de Hamming, se acepta y se refuerza con una corrección
colineal por \aLMS; (3)~si no, se prueban pares de Adalines; (4)~la capa de salida se entrena con
\aLMS. El artículo advierte que MRII puede estancarse en óptimos locales y propone como remedio la
\textbf{adición esporádica de ruido} o el reentrenamiento con distintas semillas (Sección VII-E)
\cite{widrow1990}.

\subsection{Reglas de descenso más pronunciado}\label{ssec:sd}

A diferencia de la corrección de error, estas reglas minimizan el error \emph{promediado} sobre el
conjunto de entrenamiento —típicamente el MSE— mediante descenso por gradiente, usando en la
práctica el gradiente instantáneo de un patrón a la vez.

\subsubsection{$\mu$-LMS y la superficie de MSE (Eq.~33)}

\begin{equation}
W_{k+1} = W_k + 2\mu\,\varepsilon_k\,X_k
\label{eq:mu_lms}
\end{equation}

La superficie de MSE $\xi(W) = \E[d^2] - 2P^T W + W^T R W$ es un hiperparaboloide \textbf{convexo}
con mínimo único, la solución de Wiener $W^* = R^{-1}P$ (Fig.~17). El gradiente instantáneo
$-2\varepsilon_k X_k$ es un estimador \emph{insesgado} del verdadero. Estabilidad:
$0 < \mu < 1/\mathrm{tr}[R]$. Frente a \aLMS: \muLMS~converge a la solución de mínimo MSE
(Wiener), mientras que \aLMS~converge a la solución de Wiener del conjunto \emph{normalizado} y,
para entradas no binarias, a una solución sesgada \cite{widrow1990}.

\subsubsection{Backpropagation: del Adaline sigmoidal a la red}

Para el Adaline sigmoide ($y = \tanh(s)$), minimizar $\E[\tilde\varepsilon^2]$ por descenso da la
regla (Eq.~56):

\begin{equation}
W_{k+1} = W_k + 2\mu\,\tilde{\varepsilon}_k \cdot (1 - y_k^2)\,X_k
\label{eq:backprop_single}
\end{equation}

En una red multicapa, la generalización es \bp: las ``derivadas del error cuadrático'' $\delta$ se
propagan hacia atrás. Para la capa de salida (Eq.~79) y las ocultas (Eq.~87):
\begin{align}
\delta^{(L)}_j &= \varepsilon_j \cdot (1 - y_j^2) \label{eq:delta_out}\\
\delta^{(l)} &= \bigl(W^{(l+1)T} \delta^{(l+1)}\bigr) \odot (1 - (y^{(l)})^2) \label{eq:delta_hid}
\end{align}
y la actualización, con \textbf{momentum} (Eqs.~96--97):
\begin{equation}
\Delta W_k = \eta\,\Delta W_{k-1} + (1-\eta)\,2\mu\,\delta_k X_k^T,\quad \eta \approx 0.8\text{--}0.9
\label{eq:momentum}
\end{equation}
Es crucial notar que $\delta$ \textbf{no es un error}: es una señal de gradiente que indica cómo
cambiar cada peso para reducir el error cuadrático de salida $\varepsilon^2$
\cite{widrow1990,rumelhart1986}.

\subsubsection{Madaline Rule III (MRIII) y equivalencia con Backpropagation}

\mriii~(Andes, 1988) reemplaza el signum por sigmoides y estima el gradiente de $\varepsilon^2$
respecto a $s_j$ mediante \textbf{perturbación del sumidero} (Eqs.~102--103) \cite{widrow1990}:

\begin{equation}
\frac{\partial \varepsilon^2}{\partial s_j} \approx
\frac{\varepsilon^2_{\text{perturbado}} - \varepsilon^2_{\text{base}}}{\Delta s}
\label{eq:mriii}
\end{equation}

Cuando $\Delta s \to 0$, MRIII es \textbf{matemáticamente equivalente a backpropagation}. Requiere
$1 + N_{\text{Adalines}}$ pasadas hacia adelante por patrón (frente a $1+1$ de \bp), lo que lo hace
menos eficiente en digital, pero más robusto para hardware analógico, donde no necesita conocer la
derivada exacta de la sigmoide (existió un chip comercial de MRIII) \cite{widrow1990}.

\subsection{Síntesis: taxonomía de algoritmos}\label{ssec:taxonomia}

La Figura~\ref{fig:taxonomia} reproduce la taxonomía de la Fig.~33 del artículo, que resume la
totalidad de reglas estudiadas. Las dos ramas superiores son los mecanismos de adaptación; cada una
se subdivide por arquitectura (red en capas / elemento único) y por no linealidad
(sigmoide, signum o lineal). El hallazgo central del paper —que \mriii~y \bp~coinciden— se refleja
en que ambos aparecen como hojas de la rama sigmoide del descenso más pronunciado.

\begin{figure}[htbp]\centering
\resizebox{\textwidth}{!}{%
\begin{tikzpicture}[
  every node/.style={font=\footnotesize},
  cat/.style={draw,rounded corners,fill=blue!8,align=center,inner sep=4pt,font=\small\bfseries},
  arch/.style={draw,align=center,inner sep=3pt},
  act/.style={align=center,font=\footnotesize\itshape},
  alg/.style={draw,rounded corners,fill=gray!12,align=center,inner sep=3pt,font=\footnotesize\bfseries},
  ln/.style={-}]
% root
\node (root) at (0,7) {\textbf{Reglas de aprendizaje}};
% categories
\node[cat] (sd) at (-5,5.6) {Descenso más\\pronunciado};
\node[cat] (ec) at (5,5.6) {Corrección\\de error};
\draw[ln] (root)--(sd); \draw[ln] (root)--(ec);
% architectures
\node[arch] (sdL) at (-7.6,4.2) {Red en\\capas};
\node[arch] (sdS) at (-3.0,4.2) {Elemento\\único};
\node[arch] (ecL) at (2.8,4.2) {Red en\\capas};
\node[arch] (ecS) at (7.6,4.2) {Elemento\\único};
\draw[ln] (sd)--(sdL); \draw[ln] (sd)--(sdS);
\draw[ln] (ec)--(ecL); \draw[ln] (ec)--(ecS);
% activations
\node[act] (sdLa) at (-7.6,2.9) {No lineal\\(sigmoide)};
\node[act] (sdSa) at (-3.9,2.9) {No lineal\\(sigmoide)};
\node[act] (sdSl) at (-2.0,2.9) {Lineal};
\node[act] (ecLa) at (2.8,2.9) {No lineal\\(signum)};
\node[act] (ecSa) at (6.7,2.9) {No lineal\\(signum)};
\node[act] (ecSl) at (8.6,2.9) {Lineal};
\draw[ln] (sdL)--(sdLa);
\draw[ln] (sdS)--(sdSa); \draw[ln] (sdS)--(sdSl);
\draw[ln] (ecL)--(ecLa);
\draw[ln] (ecS)--(ecSa); \draw[ln] (ecS)--(ecSl);
% leaves
\node[alg] (l1) at (-7.6,1.3) {MRIII\\Backprop};
\node[alg] (l2) at (-3.9,1.3) {MRIII\\Backprop};
\node[alg] (l3) at (-2.0,1.3) {$\mu$-LMS};
\node[alg] (l4) at (2.8,1.3) {MRI\\MRII};
\node[alg] (l5) at (6.7,1.3) {Perceptrón\\Mays};
\node[alg] (l6) at (8.6,1.3) {$\alpha$-LMS};
\draw[ln] (sdLa)--(l1); \draw[ln] (sdSa)--(l2); \draw[ln] (sdSl)--(l3);
\draw[ln] (ecLa)--(l4); \draw[ln] (ecSa)--(l5); \draw[ln] (ecSl)--(l6);
\end{tikzpicture}}
\caption{Taxonomía de las reglas de aprendizaje del artículo (reproducción de la Fig.~33 del paper).
De los ocho algoritmos, se implementan y experimentan seis: Perceptrón, \aLMS, \muLMS, MLP con
\bp, \mrii~y \mriii; MRI, Mays y el preprocesador polinomial se cubren en el marco teórico.}
\label{fig:taxonomia}
\end{figure}

% ═══════════════════════════════════════════════════════════════════
%  3. EXPERIMENTACIÓN
% ═══════════════════════════════════════════════════════════════════
\section{Experimentación}

\subsection{Diseño experimental}

Dado que el artículo es una revisión teórica, la experimentación consiste en \textbf{implementar
desde cero los algoritmos experimentables} y verificar las hipótesis H1--H4 mediante siete
experimentos (E1--E7), con validación en datasets reales adicionales. La
Tabla~\ref{tab:diseno} resume el diseño. Las reglas de Mays, MRI y el preprocesador polinomial
se analizan en el marco teórico y no tienen experimento propio.

\begin{table}[htbp]\centering\footnotesize
\begin{tabular}{@{}clllp{3.2cm}@{}}
\toprule
\textbf{Exp.} & \textbf{Tema} & \textbf{Hip.} & \textbf{Algoritmos} & \textbf{Dataset} \\
\midrule
E1 & Capacidad de Cover & H1 & Programaci\'on Lineal & Sint\'etico Monte Carlo \\
E2 & Reglas lineales, separable & H2 & Perceptr\'on, \aLMS, \muLMS & Gaussianas separables \\
E3 & Datos no separables & H2 & Perceptr\'on, \aLMS, \muLMS & Gaussianas solapadas \\
E4 & Superficies MSE & --- & Adaline 2D & Patrones aleatorios \\
E5 & XOR: MRII vs MLP & H3 & MRII, MLP-\bp & XOR (4 patrones) \\
E6 & Datasets reales & --- & MLP-\bp, \aLMS, \muLMS & Iris, Wine, MNIST-04, moons \\
E7 & MRIII $\equiv$ \bp~+ costo & H4 & MRIII, \bp & Iris (1 patr\'on) \\
\bottomrule
\end{tabular}
\caption{Diseño experimental.}\label{tab:diseno}
\end{table}

\subsection{Datasets}

Sintéticos separables (200 puntos, gaussianas $N([\pm2,\pm2], 0.35^2I)$); no separables
($N([\pm0.9,\pm0.9], 1.3^2I)$); XOR (4 patrones $\{\pm1\}^2$); Iris (Fisher, 1936; UCI: 150
flores, 4 atributos, 3 clases, \textit{split} 70/30 estratificado, normalizado a $[-1,1]$);
two-moons (300 puntos, ruido 0.20); Wine (UCI: 178 muestras, 13 atributos, binarizado); MNIST-04
(UCI digits: dígitos 0--4, $8\times8=64$ atributos, 901 instancias).

\subsection{Implementación y complejidad}

Todos los algoritmos se implementaron desde cero en \textbf{NumPy puro} (sin PyTorch, TensorFlow ni
el aprendizaje de scikit-learn, usado solo para carga de datos y métricas). La reproducibilidad se
garantiza con semilla global \texttt{SEED=1990}. El costo por presentación es $\mathcal{O}(N_w)$
para Perceptrón/LMS y para \bp~(forward $+$ backward), y $\mathcal{O}(N_w \cdot N_{\text{Adalines}})$
para MRIII, que requiere $1 + N_{\text{Adalines}}$ pasadas hacia adelante (una sin perturbar y una
por cada Adaline perturbada). Esta diferencia se cuantifica en E7.

% ═══════════════════════════════════════════════════════════════════
%  4. ANÁLISIS DE RESULTADOS
% ═══════════════════════════════════════════════════════════════════
\section{Análisis de Resultados}

\subsection{E1: Capacidad de Cover}

\figura[0.70]{fig_cover}{Probabilidad de separabilidad lineal: empírica (puntos) vs.~fórmula teórica de Cover (curvas). Líneas verticales en $C_d=N_w$ (punteada) y $C_s \approx 2N_w$ (discontinua).}{cover}

La Figura~\ref{fig:cover} muestra un ajuste estrecho entre la probabilidad empírica (80 ensayos
Monte Carlo por punto, separabilidad testeada con \texttt{linprog}) y la fórmula teórica
(Eq.~\ref{eq:cover}). Las desviaciones medias $|\text{empírica} - \text{teórica}|$ son
$0.0337$ para $N_w=2$, $0.0259$ para $N_w=5$ y $0.0105$ para $N_w=15$, decreciendo con $N_w$
como predice la teoría asintótica. La transición en $N_p/N_w = 1$ ($C_d$) y $N_p/N_w = 2$
($C_s$) se agudiza al aumentar $N_w$, confirmando \textbf{H1}.

\subsection{E2: Reglas lineales en datos separables}

\figura[0.60]{fig_datasets}{Datasets sintéticos utilizados: separable, no separable y XOR.}{datasets}

\figura{fig_fronteras}{Fronteras de decisión de los tres algoritmos lineales (color) vs.~el
hiperplano óptimo de Wiener (negro discontinuo). Ángulos respecto a $W^*$ indicados.}{fronteras}

\figura{fig_convergencia}{Convergencia de \muLMS~para distintos $\mu$. $\mu=0.12$ diverge por
exceder la cota $\mu_{\max} = 1/\mathrm{tr}[R] = 0.111$. La línea horizontal marca el MSE de
Wiener (0.015).}{convergencia}

La Figura~\ref{fig:fronteras} revela un resultado fundamental: \textbf{el Perceptrón converge
en solo 2 épocas} (confirmando el teorema), pero su frontera queda a $20.6^\circ$ del óptimo de
Wiener — el Perceptrón ``se detiene donde primero separa''. En contraste, \aLMS~y \muLMS~optimizan
el MSE, acercándose a $4.4^\circ$ y $3.5^\circ$ de $W^*$ respectivamente.

La Figura~\ref{fig:convergencia} demuestra la cota de estabilidad: con $\mu=0.12 >
1/\mathrm{tr}[R] = 1/9.02 = 0.111$, el MSE diverge exponencialmente. Los valores $\mu \in
\{0.001, 0.01, 0.05\}$ convergen al MSE de Wiener ($0.0148$), con mayor $\mu$ implicando
convergencia más rápida. Esto confirma \textbf{H2} para el caso separable.

\subsection{E3: Datos no separables}

\figura[0.72]{fig_noseparable}{Izquierda: accuracy por época en datos no separables. Derecha: evolución
de $\norm{W}$. Nótese la oscilación del Perceptrón (azul).}{noseparable}

\tabla{tab_lineales}{Resultados comparativos de algoritmos lineales en datos separables y no separables.}{lineales}

En el dataset no separable (Figura~\ref{fig:noseparable}), el \textbf{Perceptrón oscila}
significativamente: $75.7\% \pm 5.5$ en las últimas 40 épocas, con $\norm{W}$ fluctuante (el
vector ``gravita hacia cero'' como advierte el paper). \aLMS~es estable pero ruidoso
($75.0\% \pm 7.4$). \muLMS~ofrece el mejor desempeño: $79.6\% \pm 1.5$, a solo 1.4 puntos
porcentuales del óptimo de Wiener ($81.0\%$) y con la menor varianza. Esto confirma \textbf{H2}
en el caso no separable. (La regla de Mays de la Eq.~\ref{eq:mays}, con su zona muerta, atacaría
este mismo régimen y —según el paper— superaría al Perceptrón por su tendencia a alejar el vector
de pesos de cero.)

\subsection{E4: Superficies de MSE}

\figura{fig_superficies_mse}{Superficies de MSE para un Adaline de 2 pesos sin bias
(6 patrones aleatorios). Izquierda: error lineal (paraboloide convexo, Fig.~22 del paper). Centro:
error sigmoide (no cuadrático pero suave, Fig.~23). Derecha: error signum (escalonado con mesetas y
óptimos locales, Fig.~24).}{superficies}

La Figura~\ref{fig:superficies} ilustra por qué el descenso de gradiente \textbf{exige funciones
de activación diferenciables}: el error \textbf{lineal} ($\min = 0.936$) es un paraboloide
convexo (mínimo global garantizado); el \textbf{sigmoide} ($\min = 0.940$) es no cuadrático pero
suave y navegable; el \textbf{signum} ($\min = 1.333$) es escalonado, con \textbf{mesetas} donde
el gradiente es cero. Esto explica por qué algoritmos como MRII (que usan signum) requieren
mecanismos de escape como el ruido esporádico.

\subsection{E5: XOR — MRII vs MLP-Backpropagation}

\figura{fig_xor}{Izquierda: regiones de decisión de MRII 2--3--1 con ruido. Centro: MLP-\bp~2--2--1.
Derecha: curva de aprendizaje MSE del MLP.}{xor}

\tabla{tab_xor}{Tasa de éxito en XOR sobre 20 inicializaciones distintas.}{xor}

La Tabla~\ref{tab:xor} presenta el resultado más revelador del experimento. \textbf{MRII 2--2--1
sin ruido: 0/20}; \textbf{MRII con ruido esporádico ($\sigma=0.15$): 0/20} incluso con 3 Adalines
ocultas. En nuestra implementación, el ruido gaussiano no bastó para escapar de los mínimos
locales en XOR, posiblemente porque la corrección absoluta agresiva del refuerzo colineal domina
sobre el escape estocástico. Es importante señalar que el propio artículo indica que MRII
\emph{puede} resolver problemas de este tipo, y que su remedio para óptimos locales es el
reentrenamiento con distintas semillas iniciales (\textit{random restarts}) además del ruido
(Sección VII-E); nuestro resultado de 0/20 refleja la fragilidad de una única corrida por semilla,
no una imposibilidad del algoritmo. En contraste, \textbf{MLP-\bp~2--2--1: 6/20} (mediana $\approx$
40 épocas) y \textbf{MLP-\bp~2--3--1: 11/20}: una Adaline oculta adicional casi duplica la tasa de
éxito. Estos resultados confirman \textbf{H3}: XOR es resoluble por arquitecturas multicapa, pero
tanto MRII como \bp~están sujetos a óptimos locales, y la superficie sigmoide (E4) es
inherentemente más navegable que la del signum.

\subsection{E6: MLP en datasets reales}

\figura{fig_iris_mlp}{MLP 4--8--3 en Iris. Izquierda: curva MSE. Centro: accuracy por época.
Derecha: matriz de confusión (91.1\% test).}{iris_mlp}

\figura[0.55]{fig_moons}{Regiones de decisión del MLP 2--8--1 en two-moons (93.3\% test).}{moons}

\figura[0.60]{fig_wine}{Convergencia de clasificadores lineales en Wine. Wiener y \muLMS~alcanzan
98.1\% test.}{wine}

\figura[0.58]{fig_mnist}{MLP 64--16--5 en MNIST dígitos 0--4 (98.5\% test).}{mnist}

\tabla{tab_iris}{Resultados del MLP 4--8--3 en Iris para la grilla $\mu \times \eta$.}{iris}

\tabla{tab_relevancia}{Relevancia de variables en Iris medida como $\norm{W_1}_2$ por atributo.}{relevancia}

\textbf{Iris (Figura~\ref{fig:iris_mlp}, Tablas~\ref{tab:iris}--\ref{tab:relevancia}):} Setosa es
linealmente separable del resto (el Perceptrón converge en 2 épocas, 100\% test), pero Versicolor
vs.~Virginica \textbf{no} lo es (techo $\approx 83\%$ lineal). El MLP 4--8--3 alcanza
consistentemente \textbf{91.1\% en test} para todas las configuraciones de $\mu$ y $\eta$, con
convergencia más rápida para $\mu=0.05$. El momentum ($\eta=0.9$) no aporta beneficio en este
problema pequeño y bien condicionado, consistente con el paper: ``el momentum por sí solo suele ser
de poco valor'' \cite{widrow1990}. El \textbf{criterio de selección del modelo final} favorece
$\mu=0.05,\ \eta=0$ por su convergencia más rápida a igual accuracy. El análisis de relevancia
(Tabla~\ref{tab:relevancia}) revela que los pétalos concentran el \textbf{74.4\%} del peso total de
la primera capa (\textit{petal width}: 48.0\%, \textit{petal length}: 26.4\%), consistente con el
conocimiento botánico de que son los discriminantes más informativos \cite{fisher1936}.

\textbf{Two-moons, Wine y MNIST:} el MLP 2--8--1 alcanza \textbf{93.3\%} en two-moons trazando una
frontera no lineal; los clasificadores lineales alcanzan \textbf{98.1\%} en Wine (la clase 0 es
mayoritariamente separable en el espacio de 13 atributos); y el MLP 64--16--5 alcanza
\textbf{98.5\%} en MNIST-04, demostrando que una arquitectura simple con \bp~maneja datos de
mediana dimensionalidad con alta precisión.

\subsection{E7: MRIII $\equiv$ Backpropagation}

\figura{fig_mriii_bp}{Izquierda: error relativo $\norm{\nabla_{\text{MRIII}} -
\nabla_{\text{BP}}} / \norm{\nabla_{\text{BP}}}$ vs.~$\Delta s$ (log-log, eje $x$ invertido);
mínimo $5.05 \times 10^{-8}$ en $\Delta s = 10^{-8}$. Derecha: costo computacional por presentación
(ms).}{mriii_bp}

\tabla{tab_mriii}{Costo computacional de \bp~vs.~MRIII para redes 4--H--3 (tiempos dependientes del hardware).}{mriii}

La Figura~\ref{fig:mriii_bp}(a) confirma \textbf{H4} con precisión numérica: el gradiente estimado
por perturbación (MRIII) converge al gradiente analítico de \bp~con error relativo mínimo de
$\mathbf{5.05 \times 10^{-8}}$ en $\Delta s = 10^{-8}$. El error decrece linealmente con $\Delta s$
hasta que la precisión de punto flotante de 64 bits ($\approx 10^{-16}$) provoca un rebote para
$\Delta s < 10^{-8}$.

La Figura~\ref{fig:mriii_bp}(b) y la Tabla~\ref{tab:mriii} cuantifican el costo de esa
equivalencia. La razón MRIII/\bp~\textbf{crece aproximadamente como
$\mathcal{O}(N_{\text{Adalines}})$}: pasa de unas pocas veces para redes pequeñas ($H=4$) a dos
órdenes de magnitud para las grandes ($H=256$). Los valores absolutos de la tabla son dependientes
del hardware y de la implementación, pero la tendencia lineal en $N_{\text{Adalines}}$ es robusta.
Esto valida que \bp~es la opción óptima para implementación digital, mientras que MRIII tiene
sentido en hardware analógico, donde la perturbación física del sumidero es más simple que la
retropropagación de señales de error.

\subsection{Tabla resumen}

\tabla{tab_resumen}{Resumen consolidado de todos los experimentos.}{resumen}

% ═══════════════════════════════════════════════════════════════════
%  5. CONCLUSIONES
% ═══════════════════════════════════════════════════════════════════
\section{Conclusiones}

\subsection{Aporte al conocimiento}

El artículo de Widrow \& Lehr (1990) realiza un aporte fundamental al \textbf{unificar 30 años de
algoritmos de aprendizaje supervisado} para RNA feedforward bajo un principio común: la mínima
perturbación. Los experimentos verifican empíricamente esta unificación en tres planos:

\begin{enumerate}[leftmargin=*]
    \item \textbf{El linaje Perceptrón $\to$ Adaline $\to$ Madaline $\to$ Backpropagation} no es
    una colección arbitraria, sino una progresión gobernada por el mismo principio geométrico:
    $\Delta W_k$ es siempre colineal con $X_k$, minimizando $\norm{\Delta W_k}$ para la corrección
    deseada. La taxonomía de la Fig.~\ref{fig:taxonomia} organiza toda esta familia.
    \item \textbf{El trade-off capacidad--generalización} ($C_s \approx 2N_w$, $N_p \gg N_w/N_y$)
    establece límites que siguen vigentes en la era del \textit{deep learning}: una red
    sobreparametrizada memoriza; una bien dimensionada generaliza.
    \item \textbf{El valor pedagógico e histórico}: las ecuaciones de \bp~(1986) ya estaban
    implícitas en Werbos (1974), y la equivalencia MRIII $\equiv$ \bp~muestra que el descubrimiento
    científico ocurre por múltiples caminos independientes.
\end{enumerate}

\subsection{Evaluación de las hipótesis}

\begin{description}
    \item[H1 -- Capacidad de Cover:] \textbf{Confirmada.} El Monte Carlo (E1) reproduce la fórmula
    teórica con desviación media $\leq 0.034$.
    \item[H2 -- Convergencia del Perceptrón:] \textbf{Confirmada.} Converge en 2 épocas en datos
    separables (E2), pero oscila $\pm5.5$ pts en no separables (E3), mientras \muLMS~mantiene
    estabilidad ($\pm1.5$ pts, a 1.4 pts de Wiener).
    \item[H3 -- Poder no lineal y óptimos locales:] \textbf{Confirmada con matices.} XOR es
    resoluble solo por arquitecturas multicapa, pero MRII (0/20 por semilla única) resultó más
    vulnerable a óptimos locales que MLP-\bp~(6/20 y 11/20). El artículo prevé este comportamiento
    y recomienda \textit{random restarts} y ruido como remedio.
    \item[H4 -- Equivalencia MRIII $\equiv$ Backpropagation:] \textbf{Confirmada.} El error relativo
    entre gradientes es $5.05 \times 10^{-8}$, y el costo de MRIII escala como
    $\mathcal{O}(N_{\text{Adalines}})$ (de pocas veces a dos órdenes de magnitud según el tamaño de
    la red y el hardware).
\end{description}

\subsection{Limitaciones}

Los datasets son de baja dimensionalidad ($\leq 64$ atributos); no se evaluó la escalabilidad a
alta dimensión. La implementación de MRII no logró resolver XOR con una sola corrida por semilla;
una estrategia de \textit{random restarts} o recocido simulado —prevista por el paper— mejoraría el
resultado. Del lado teórico, MRI, las reglas de Mays y el preprocesador polinomial se cubren
conceptualmente pero sin experimento propio. No se implementaron variantes avanzadas de \bp
(\textit{quickprop}, \textit{delta-bar-delta}, \textit{backprop through time}) ni métodos modernos
(SVM, \textit{transformers}) como línea base.

\subsection{Reflexión personal}

Treinta y cinco años después, el artículo permanece vigente. El principio de mínima perturbación es,
en esencia, el mismo que gobierna el \textit{stochastic gradient descent} que entrena los modelos
de \textit{deep learning} actuales: \muLMS~generalizado a arquitecturas profundas. El momentum
(Eqs.~96--97) sigue siendo estándar en optimizadores como Adam, y la perturbación de MRIII anticipa
los métodos de optimización de orden cero (\textit{zeroth-order optimization}) usados hoy cuando el
gradiente no está disponible. La gran lección es que los fundamentos importan: comprender la
geometría del aprendizaje —el hiperplano separador, el paraboloide de MSE, las mesetas del error
signum— es esencial para diseñar, depurar y mejorar cualquier sistema de aprendizaje automático.

% ═══════════════════════════════════════════════════════════════════
%  REFERENCIAS
% ═══════════════════════════════════════════════════════════════════
\begin{thebibliography}{99}

\bibitem{widrow1990}
B.~Widrow and M.~A.~Lehr,
``30 years of adaptive neural networks: Perceptron, Madaline, and backpropagation,''
\textit{Proceedings of the IEEE}, vol.~78, no.~9, pp.~1415--1442, 1990.

\bibitem{rosenblatt1958}
F.~Rosenblatt,
``The perceptron: A probabilistic model for information storage and organization in the brain,''
\textit{Psychological Review}, vol.~65, no.~6, pp.~386--408, 1958.

\bibitem{widrow1960}
B.~Widrow and M.~E.~Hoff,
``Adaptive switching circuits,''
in \textit{IRE WESCON Convention Record}, vol.~4, 1960, pp.~96--104.

\bibitem{cover1965}
T.~M.~Cover,
``Geometrical and statistical properties of systems of linear inequalities with applications
in pattern recognition,''
\textit{IEEE Transactions on Electronic Computers}, vol.~EC-14, no.~3, pp.~326--334, 1965.

\bibitem{rumelhart1986}
D.~E.~Rumelhart, G.~E.~Hinton, and R.~J.~Williams,
``Learning representations by back-propagating errors,''
\textit{Nature}, vol.~323, pp.~533--536, 1986.

\bibitem{werbos1974}
P.~J.~Werbos,
``Beyond regression: New tools for prediction and analysis in the behavioral sciences,''
Ph.D.~dissertation, Harvard University, Cambridge, MA, 1974.

\bibitem{mays1963}
C.~H.~Mays,
``Adaptive threshold logic,''
Ph.D.~dissertation, Tech.~Rep.~1557-1, Stanford Electron.~Labs., Stanford, CA, 1963.

\bibitem{widrow1962}
B.~Widrow,
``Generalization and information storage in networks of adaline `neurons','' in
\textit{Self-Organizing Systems 1962}, Washington, DC: Spartan Books, 1962, pp.~435--461.

\bibitem{specht1967}
D.~F.~Specht,
``Generation of polynomial discriminant functions for pattern recognition,''
\textit{IEEE Transactions on Electronic Computers}, vol.~EC-16, pp.~308--319, 1967.

\bibitem{fisher1936}
R.~A.~Fisher,
``The use of multiple measurements in taxonomic problems,''
\textit{Annals of Eugenics}, vol.~7, no.~2, pp.~179--188, 1936.

\bibitem{widrow1987}
B.~Widrow, R.~G.~Winter, and R.~A.~Baxter,
``Layered neural nets for pattern recognition,''
\textit{IEEE Transactions on Acoustics, Speech, and Signal Processing}, vol.~36, no.~7,
pp.~1109--1118, 1988.

\bibitem{nguyen1989}
D.~H.~Nguyen and B.~Widrow,
``Neural networks for self-learning control systems,''
\textit{IEEE Control Systems Magazine}, vol.~10, no.~3, pp.~18--23, 1990.

\bibitem{andes1990}
D.~Andes, B.~Widrow, M.~Lehr, and E.~Wan,
``MRIII: A robust algorithm for training analog neural networks,''
\textit{Proc.~Int.~Joint Conf.~on Neural Networks}, vol.~1, pp.~533--536, 1990.

\bibitem{bishop2006}
C.~M.~Bishop,
\textit{Pattern Recognition and Machine Learning}.
New York: Springer, 2006.

\end{thebibliography}

\end{document}
